% =============================================================================
%  Notas del Profesor — Sesión 17: Taller 1 — Resolución de problemas
%  Programación Avanzada — Universidad Panamericana
%  Compilar: pdflatex sesion17_taller_1_profesor.tex
% =============================================================================
\documentclass[12pt, oneside]{article}
\usepackage[profesor]{progavanzada}

\begin{document}

\portada{Sesión 17}{Taller 1 — Resolución de problemas}{2026}

\tableofcontents
\newpage

% =============================================================================
\section{Objetivos de la sesión}
% =============================================================================

Al finalizar la sesión el alumno será capaz de:

\begin{enumerate}
  \item Leer un enunciado de concurso y extraer la condición matemática central.
  \item Convertir una condición con promedios en aritmética entera
        (evitar flotantes innecesarios).
  \item Implementar la reversión de un entero con operaciones de dígitos.
  \item Enviar una solución a un juez en línea e interpretar el veredicto.
\end{enumerate}

\begin{notaprofesor}[Duración estimada]
  90 minutos.  Distribución sugerida:\\
  10 min — explicación de jueces en línea + dinámica del taller \\
  35 min — Problema 1 (Paradoja de los promedios) \\
  35 min — Problema 2 (Ayudando a mi hermanito) \\
  10 min — cierre: discusión de soluciones y casos esquina \\[4pt]
  Los alumnos trabajarán en el cuaderno Jupyter
  (\texttt{sesiones/sesion17\_taller\_1.ipynb}) y enviarán al juez la
  solución en C++ al terminar.
\end{notaprofesor}

% =============================================================================
\section{Dinámica del taller}
% =============================================================================

\begin{notaprofesor}[Cómo arrancar]
  \textbf{Preguntas de sondeo (5 min):}
  \begin{itemize}
    \item ``¿Alguien ha enviado código a un juez en línea antes?''
          — Esperar manos; si hay, pedirles que cuenten su experiencia.
    \item ``Si saco a la persona con el promedio más bajo del salón y la
          mando a otro salón, ¿qué pasa con el promedio de este salón?''
          — Respuesta esperada: sube.  ``¿Y con el del otro?''
          — Depende de si está por debajo del promedio de ese salón.
          Este es exactamente el problema 1.
  \end{itemize}
\end{notaprofesor}

La sesión es \textbf{individual}.  Los alumnos tienen acceso al cuaderno
Jupyter para probar su código localmente antes de enviar.  Deben obtener
\texttt{AC} en ambos problemas.

\begin{notaprofesor}[Cuenta omegaUp]
  Asegúrate de que todos los alumnos tengan cuenta en omegaUp antes de
  la sesión.  El problema 2 está en:
  \url{https://www.omegaup.com/arena/problem/La-tarea-de-mi-hermanito/}
\end{notaprofesor}

\newpage

% =============================================================================
\section{Problema 1: Paradoja de los promedios}
% =============================================================================

\subsection{Enunciado completo}

Un chiste conocido dice que si un mal estudiante de Cómputo se cambia a
Economía, sube el promedio de inteligencia en ambas facultades.

Dado el IQ de todos los estudiantes de dos facultades, determina cuántos
estudiantes de Cómputo podrían protagonizar el chiste (es decir, cuántos
harían que \textbf{ambos} promedios suban al cambiar de facultad).

\textbf{Entrada:} primera línea con $T$ casos.  Cada caso precedido por
línea en blanco.  Cada caso: primera línea con $n$ y $m$ (al menos
2 estudiantes de Cómputo); las siguientes líneas contienen $n + m$
enteros positivos (primeros $n$ son Cómputo, los $m$ restantes Economía).

\textbf{Salida:} un entero por caso.

\textbf{Muestra:} $n=5$, $m=5$, Cómputo: 100 101 102 103 104,
Economía: 98 100 102 99 101.  Salida: 1.

\subsection{Análisis de la solución}

Sea $\bar{c} = \text{sum\_cs}/n$ el promedio de Cómputo y
$\bar{e} = \text{sum\_ec}/m$ el promedio de Economía.

Un estudiante con IQ $x$ cumple el chiste si y solo si:

\medskip
\textbf{Condición 1 — sube el promedio de Cómputo al salir:}
\[
  \frac{\text{sum\_cs} - x}{n - 1} > \frac{\text{sum\_cs}}{n}
  \;\Longleftrightarrow\;
  x \cdot n < \text{sum\_cs}
  \;\Longleftrightarrow\;
  x < \bar{c}
\]

\textbf{Condición 2 — sube el promedio de Economía al entrar:}
\[
  \frac{\text{sum\_ec} + x}{m + 1} > \frac{\text{sum\_ec}}{m}
  \;\Longleftrightarrow\;
  x \cdot m > \text{sum\_ec}
  \;\Longleftrightarrow\;
  x > \bar{e}
\]

Condición combinada: $\bar{e} < x < \bar{c}$.

\begin{notaprofesor}[Clave de la solución]
  Si $\bar{c} \leq \bar{e}$, ningún estudiante cumple (el rango está vacío).
  La verificación entera $x \cdot n < \text{sum\_cs}$ y
  $x \cdot m > \text{sum\_ec}$ evita comparar doubles.
  Con IQs hasta $10^4$ y hasta $10^4$ estudiantes, el producto cabe en
  \texttt{long long}.
\end{notaprofesor}

\subsection{Verificación con la muestra}

$\text{sum\_cs} = 510$, $\bar{c} = 102$. \quad
$\text{sum\_ec} = 500$, $\bar{e} = 100$.

\begin{center}
\begin{tabular}{cccc}
  \toprule
  $x$ & $x \cdot 5 < 510$? & $x \cdot 5 > 500$? & ¿Cuenta? \\
  \midrule
  100 & 500 < 510 \checkmark & 500 > 500 \texttimes & No \\
  101 & 505 < 510 \checkmark & 505 > 500 \checkmark & \textbf{Sí} \\
  102 & 510 < 510 \texttimes & — & No \\
  103 & 515 < 510 \texttimes & — & No \\
  104 & 520 < 510 \texttimes & — & No \\
  \bottomrule
\end{tabular}
\end{center}

Resultado: \textbf{1}. \checkmark

\subsection{Solución en C++}

\begin{verbatim}
#include <iostream>
#include <vector>
using namespace std;

int main() {
    int T;
    cin >> T;
    while (T--) {
        int n, m;
        cin >> n >> m;
        vector<long long> cs(n), ec(m);
        for (auto& x : cs) cin >> x;
        for (auto& x : ec) cin >> x;

        long long sum_cs = 0, sum_ec = 0;
        for (auto x : cs) sum_cs += x;
        for (auto x : ec) sum_ec += x;

        int cnt = 0;
        for (auto x : cs)
            if (x * n < sum_cs && x * m > sum_ec)
                cnt++;

        cout << cnt << "\n";
    }
}
\end{verbatim}

\begin{notaprofesor}[FAQs Problema 1]
  \textbf{¿Por qué no basta con calcular los promedios como \texttt{double}?}\\
  Con IQs enteros y promedios exactos la diferencia puede ser 0.5 o menos.
  La comparación $x < \bar{c}$ con \texttt{double} puede dar falsos positivos
  o negativos por error de representación.  Con enteros es exacto.

  \medskip
  \textbf{¿Qué pasa si $n = 2$ y el estudiante que sale es el único que lo
  haría?}\\
  No hay caso especial: la condición $x \cdot n < \text{sum\_cs}$ sigue
  siendo válida para $n \geq 2$.  Al salir quedan $n - 1 \geq 1$ estudiantes,
  que es un promedio bien definido.

  \medskip
  \textbf{¿Por qué \texttt{long long}?}\\
  Si los IQs llegan a $10^4$ y $n$ llega a $10^4$, el producto
  $x \cdot n$ puede ser $10^8$, que cabe en \texttt{int} (máx.\ $\approx
  2 \times 10^9$), pero sum\_cs con $10^4$ elementos de valor $10^4$ es
  $10^8$, que también cabe.  Usar \texttt{long long} es buena práctica
  defensiva aunque no sea estrictamente necesario aquí.
\end{notaprofesor}

\newpage

% =============================================================================
\section{Problema 2: Ayudando a mi hermanito}
% =============================================================================

\subsection{Enunciado}

Se dan dos números de tres dígitos $A$ y $B$ sin ceros.  Se invierten ambos
y se devuelve el mayor de los dos \textit{ya invertidos}, escrito en reversa.

\textbf{Ejemplos:}
\begin{center}
\begin{tabular}{llll}
  \toprule
  $A$ & $B$ & reversa($A$), reversa($B$) & Salida \\
  \midrule
  735 & 892 & 537, 298 & 537 \\
  331 & 341 & 133, 143 & 143 \\
  \bottomrule
\end{tabular}
\end{center}

\subsection{Análisis de la solución}

Para un número $n$ de tres dígitos con dígitos $d_1 d_2 d_3$ (centenas,
decenas, unidades):
\[
  \text{rev}(n) = (n \bmod 10) \times 100
               + \bigl((n / 10) \bmod 10\bigr) \times 10
               + (n / 100)
\]

Como $A$ y $B$ no contienen ceros, la reversa de un número de tres dígitos
siempre es también un número de tres dígitos (no hay ceros que desaparezcan
al principio).

\begin{notaprofesor}[Clave de la solución]
  El problema es directo una vez que se identifica la fórmula de reversión.
  La restricción ``sin ceros'' elimina el único caso especial (cero inicial
  en la reversa que haría que la salida tuviera menos de tres dígitos).
  No se necesita \texttt{string}.
\end{notaprofesor}

\subsection{Solución en C++}

\begin{verbatim}
#include <iostream>
using namespace std;

int rev(int n) {
    return (n % 10) * 100 + ((n / 10) % 10) * 10 + n / 100;
}

int main() {
    int A, B;
    cin >> A >> B;
    int rA = rev(A), rB = rev(B);
    cout << (rA > rB ? rA : rB) << "\n";
}
\end{verbatim}

\begin{notaprofesor}[FAQs Problema 2]
  \textbf{¿Por qué no funciona con ceros?}\\
  Si $A = 300$, su reversa es 003, que como entero es 3 (dos dígitos menos).
  La salida sería 3 y no 003.  El enunciado garantiza que esto no ocurre.

  \medskip
  \textbf{Un alumno usa \texttt{string} para invertir.  ¿Es correcto?}\\
  Sí, también funciona.  La solución con operaciones aritméticas es más
  eficiente y muestra mejor comprensión de la estructura de dígitos.

  \medskip
  \textbf{¿Es posible que las reversas sean iguales?}\\
  El enunciado dice que $A \neq B$.  Pero podría ocurrir que
  $\text{rev}(A) = \text{rev}(B)$ (e.g.\ $A = 123$, $B = 321$ →
  ambas reversas son 321 y 123 respectivamente — en realidad no).
  Más: $A = 129$, $B = 921$ → rev(A) = 921, rev(B) = 129, distintas.
  En la práctica es muy difícil que coincidan dado que los originales son
  distintos.  El enunciado no especifica qué hacer en ese caso; podemos
  asumir que no ocurre.
\end{notaprofesor}

\newpage

% =============================================================================
\section{Materiales y recursos}
% =============================================================================

\begin{itemize}
  \item \textbf{Cuaderno:} \texttt{sesiones/sesion17\_taller\_1.ipynb}
  \item \textbf{Notas alumno:} \texttt{notas\_alumno/pdf/sesion17\_taller\_1.pdf}
  \item \textbf{Problema 2 en omegaUp:}
        \url{https://www.omegaup.com/arena/problem/La-tarea-de-mi-hermanito/}
\end{itemize}

\end{document}
